% Lernkarten -- Analysis (Modul 61211), Lektion 1
% Quelle: eigene OneNote-Notizen (Fernuni Notizbuch > 61211 Analysis)
% Format: \lernkarte{Vorderseite (roter Titel)}{Rueckseite (weisser Inhalt)}
%
% --- OneNote-Seite 1.1 "Reelle Zahlen" -----------------------------------

\lernkarte{Geordnete Menge}{%
  \begin{align*}
    & x \le x \\
    & x \le y \,\wedge\, y \le x \;\Rightarrow\; x = y \\
    & x \le y \,\wedge\, y \le z \;\Rightarrow\; x \le z
  \end{align*}
  \[ \Rightarrow \text{Axiom von } \mathbb{R} \]
}

\lernkarte{linear geordnete Menge}{%
  \[ x \le y \;\vee\; y \le x \]
  nach dieser Regel können je zwei Elemente verglichen werden.
}

\lernkarte{Regeln der Vergleichung}{%
  \[ 0 \le a \le b \,\wedge\, 0 \le c \le d \;\Rightarrow\; ac \le bd \]
}

\lernkarte{Cauchy-Schwarz Ungleichung}{%
  \[ \left|\sum_{k=1}^{n} x_k y_k\right|
     \le \sqrt{\sum_{k=1}^{n} x_k^2}\;\sqrt{\sum_{k=1}^{n} y_k^2} \]
  \[ \sqrt{\sum_{k=1}^{n} (x_k+y_k)^2}
     \le \sqrt{\sum_{k=1}^{n} x_k^2} + \sqrt{\sum_{k=1}^{n} y_k^2} \]
}

\lernkarte{Supremum \& Infimum}{%
  \begin{align*}
    \sup M &= S \le S' \in \text{obere Schranke} \\
    \inf M &= S \ge S' \in \text{untere Schranke}
  \end{align*}
  \[ \text{Maximum} \;\Leftrightarrow\; \sup M \in M \qquad
     \text{Minimum} \;\Leftrightarrow\; \inf M \in M \]
}

\lernkarte{beschränkte Funktion}{%
  \[ \|f\|_\infty := \sup\{\, \|f(x)\| \;\mid\; x \in M \,\} \]
}

\lernkarte{Supremumsnorm}{%
  \[ \|f\|_\infty < \infty \;\Rightarrow\; f \in B(M) \]
}

% --- OneNote-Seite 1.2 "Konvergenz" --------------------------------------

\lernkarte{zweite Dreiecksungleichung}{%
  \[ \bigl|\,|x|-|y|\,\bigr| \;\le\; |x+y| \]
}

\lernkarte{Umgebung}{%
  \[ U_\varepsilon(a) := \{\, x \in \mathbb{R} \;\mid\; d(x,a) < \varepsilon \,\} \]
  \[ U_\varepsilon(a) \subseteq U \;\Rightarrow\; U \text{ ist Umgebung von } a \]
}

\lernkarte{Hausdorffeigenschaft}{%
  \[ a \neq b \;\Rightarrow\; U(a) \cap U(b) = \emptyset \]
}

% PRÜFEN: Notiz schreibt wörtlich "endlich viele" (unten transkribiert).
% Mathematisch müsste es "fast alle" heißen (so auch in Seite 1.5). Bitte
% in OneNote prüfen -- ich korrigiere es nicht selbst.
\lernkarte{konvergente Folge}{%
  $f$ konvergiert gegen $a :\Leftrightarrow$ endlich viele Folgenglieder
  liegen in $U(a)$, $U(a)$ beliebig.
}

\lernkarte{Menge der konvergenten Folgen}{%
  Bezeichnung: $c$
  \[ c \subsetneq W \quad (\text{Menge aller Folgen}) \]
}

\lernkarte{Konvergenzkriterien}{%
  \begin{itemize}[topsep=2pt,itemsep=3pt,leftmargin=*]
    \item \textbf{Nullfolge:}
      $\lim f = 0 \,\wedge\, h \in \ell^\infty \;\Rightarrow\; \lim hf = 0$
    \item \textbf{Einschnürungssatz:}
      $\lim f = \lim g \,\wedge\, f_k \le h_k \le g_k \ \forall\,\text{fast } k\in\mathbb{N}
       \;\Rightarrow\; \lim f = \lim h = \lim g$
    \item \textbf{$\varepsilon$-$n_0$-Kriterium:}
      $\forall \varepsilon>0 \ \exists n_0\in\mathbb{N} \ \forall k\in\mathbb{N}:
       k \ge n_0 \Rightarrow |a_k - a| < \varepsilon$
    \item \textbf{Monotoniekriterium:} monoton wachsend $\wedge$ beschränkt
      $\;\Rightarrow\; \lim f = \sup\{\, a_k \mid k\in\mathbb{N} \,\}$
    \item \textbf{Divergenz:} eine Teilfolge divergiert, oder zwei Teilfolgen
      haben andere Grenzwerte
  \end{itemize}
}

\lernkarte{Bolzano-Weierstraß}{%
  \begin{enumerate}[topsep=2pt,itemsep=2pt,leftmargin=*]
    \item Jede Folge enthält eine monotone Teilfolge.
    \item Jede beschränkte Folge enthält eine konvergente Teilfolge.
  \end{enumerate}
}

\lernkarte{bestimmte Divergenz (Kriterium)}{%
  Für jedes $n_0 \in \mathbb{N}$ liegen fast alle $a_k$ über $n_0$.
}

\lernkarte{Cauchykriterium}{%
  Jede Cauchyfolge konvergiert.
  \begin{align*}
    &\forall \varepsilon>0 \ \exists n_0\in\mathbb{N} \ \forall k\in\mathbb{N}:
      \ k \ge n_0 \Rightarrow |a_k - a_{n_0}| < \varepsilon \\
    \Leftrightarrow\ &\forall \varepsilon>0 \ \exists n_0\in\mathbb{N} \ \forall k,l\in\mathbb{N}:
      \ k \ge n_0 \wedge l \ge n_0 \Rightarrow |a_k - a_l| < \varepsilon
  \end{align*}
}

% --- OneNote-Seite 1.3 "R^n reeller Vektorraum" --------------------------
% Hinweis: Der obere Block "Beispiele reelle Vektorräume" ist in den Notizen
% durchgestrichen und daher hier ausgelassen.

\lernkarte{$B(M)$}{%
  Beschränkte Funktionen. \\
  Unterraum von $\mathrm{Abb}(M,\mathbb{R})$.
}

\lernkarte{$C(M)$, $C^1(M)$}{%
  \begin{itemize}[topsep=2pt,itemsep=1pt,leftmargin=*]
    \item $C(M) = $ stetige Funktionen
    \item $C^1(M) = $ differenzierbar und Ableitung stetig
    \item $C^1(M)$ ist Unterraum von $C(M)$
  \end{itemize}
}

\lernkarte{Folgenräume $W, \ell^\infty, c, c_0$}{%
  Unterräume:
  \[ W \supseteq \ell^\infty \supseteq c \supseteq c_0 \]
}

% --- OneNote-Seite 1.4 "R^n normierter Raum" -----------------------------

\lernkarte{Norm}{%
  \begin{align*}
    \text{(i)} \ \ & \|x\| \ge 0 \ \wedge \ \big(\|x\| = 0 \Leftrightarrow x = 0\big) \\
    \text{(ii)} \ \ & \|\alpha x\| = |\alpha|\,\|x\| \\
    \text{(iii)} \ \ & \|x+y\| \le \|x\| + \|y\|
  \end{align*}
}

\lernkarte{Supremumsnorm (Raum)}{%
  \[ \big(\, B(M),\ \|\cdot\|_\infty \,\big) \]
}

\lernkarte{Euklidische Norm}{%
  \[ \|x\|_2 := \sqrt{\sum_{k=1}^{n} x_k^2} \]
}

\lernkarte{Maximumnorm}{%
  \[ \|x\|_\infty := \max\{\, |x_k| \;\mid\; 1 \le k \le n \,\} \]
}

% --- OneNote-Seite 1.5 "Konvergenz in R^n" -------------------------------

\lernkarte{Umgebung (in $\mathbb{R}^n$)}{%
  \[ U_\varepsilon(a) := \{\, x \in X \;\mid\; d(x,a) < \varepsilon \,\} \]
  Der Abstand $d$ wird aus der Norm $\|\cdot\|$ erzeugt.
  \[ U_\varepsilon(a) \subseteq U \;\Rightarrow\; U \text{ ist Umgebung von } a \]
}

\lernkarte{konvergente Folge (in $\mathbb{R}^n$)}{%
  In jeder Umgebung von $a$ liegen fast alle Folgenglieder
  $\;\Rightarrow\; f$ konvergiert gegen $a$.
}

\lernkarte{Konvergenzkriterium (in $\mathbb{R}^n$)}{%
  $\varepsilon$-$n_0$-Kriterium:
  \[ \forall \varepsilon>0 \ \exists n_0\in\mathbb{N} \ \forall k\in\mathbb{N}:
     \ k \ge n_0 \Rightarrow \|y_k - a\| < \varepsilon \]
  \[ \Leftrightarrow\ \lim_{k\to\infty} \|x_k - a\| = 0 \]
}

\lernkarte{Bolzano-Weierstraß (in $\mathbb{R}^n$)}{%
  \[ \|x_k\|_\infty \le S \;\Rightarrow\; \text{es existiert eine konvergente Teilfolge} \]
}

\lernkarte{Äquivalenz der Normen}{%
  \[ \|x\| \le a\,\|x\|_\infty \qquad \|x\|_\infty \le b\,\|x\| \]
  $U$ ist Umgebung von $a$ bzgl.\ $\|\cdot\|$
  $\Leftrightarrow$ bzgl.\ $\|\cdot\|_*$. \\[0.3em]
  Konvergenz in $(\mathbb{R}^n, \|\cdot\|) \Leftrightarrow (\mathbb{R}^n, \|\cdot\|_*)
  \Leftrightarrow$ komponentenweise in $(\mathbb{R}, |\cdot|)$.
}

\lernkarte{Cauchyfolge (in $\mathbb{R}^n$)}{%
  \begin{align*}
    &\forall \varepsilon>0 \ \exists n_0\in\mathbb{N} \ \forall k\in\mathbb{N}:
      \ k \ge n_0 \Rightarrow \|x_k - x_{n_0}\| < \varepsilon \\
    \Leftrightarrow\ &\forall \varepsilon>0 \ \exists n_0\in\mathbb{N} \ \forall k,l\in\mathbb{N}:
      \ k \ge n_0 \wedge l \ge n_0 \Rightarrow \|x_k - x_l\| < \varepsilon
  \end{align*}
}

\lernkarte{Banachraum}{%
  Vollständiger normierter Raum (jede Cauchyfolge konvergiert). \\
  $\mathbb{R}^n$ ist ein Banachraum.
}
